\section{Exact pairing concentration and its quantitative limitation}
\label{sec:path-seven-pure-pair-notes}

Critical seven-cut profiles have disjoint four-terminal pairing supports on their two shores.
Since a four-set has only three pairings, at least one shore is concentrated on a single pairing.
The following observation excludes exact concentration in the relevant finite graphs.
It does not give a uniform gap as rank grows.

\begin{lemma}[Alternative pairings in a bridgeless subcubic graph]
\label{lem:path-seven-alternative-pairing}
Let $X$ be a connected bridgeless subcubic graph with distinct degree-two vertices $a,b,c,d$.
Then
\[
 F_{\{a,b,c,d\}}(X)-f_{ab|cd}(X)\ge1.
\]
In particular, a nonzero four-terminal forest sector cannot be supported on only one pairing.
\end{lemma}

\begin{proof}
Add a new edge $ab$, retaining parallel edges if necessary.
The resulting graph remains subcubic and bridgeless.
A connected bridgeless subcubic graph has no cut vertex.
Each component after deleting a cut vertex would need at least two incident edges at that vertex, since a unique edge would be a bridge.
Two components would therefore force degree at least four.
Thus the resulting graph is two-connected.

In a two-connected graph, every edge belongs to a simple path between any prescribed distinct vertices.
For completeness, subdivide the edge with a new vertex and apply the two-fan form of Menger's theorem from that vertex to the prescribed pair.
Suppressing the subdivision yields the required path through the edge.
Apply this fact to the new edge $ab$ and terminals $c,d$.
Deleting $ab$ from that path gives a four-terminal forest with pairing $ac|bd$ or $ad|bc$.
This proves the stated lower bound.
\end{proof}

\begin{lemma}[Small poles have no internal bridge]
\label{lem:path-seven-poles-bridgeless}
Let $2\le k\le7$, and let a normalized $k$-pole $X$ lie in a cubic completion $G$.
If $G$ has no cyclic cut of size at most four, then $X$ is bridgeless.
\end{lemma}

\begin{proof}
Suppose $uv$ is a bridge of $X$.
Each shore of that bridge contains a cycle.
Indeed, all its vertices have degree at least two inside the shore except possibly its bridge endpoint.
A nontrivial tree needs two leaves, while a singleton shore would give a degree-one vertex in $X$.

One shore contains at most $\lfloor k/2\rfloor$ ports.
Its boundary in $G$ consists of the bridge and at most one further incidence for each of those ports.
Hence its boundary has size at most $1+\lfloor k/2\rfloor\le4$.
Both shores retain their cycles in $G$, giving the excluded cyclic cut.
\end{proof}

Together, these lemmas give an alternative-pairing mass of at least $2^{-h}$ in every four-sector of a rank-$h$ pole in this class.
That lower bound tends to zero.
It therefore does not exclude a limiting profile concentrated on one pairing.
A rank-independent lower bound, or a bound strong enough to absorb the other error terms, would be needed for that application.
No such bound is proved in this note.
